\section*{Capítulo \ref{ch:b1:field-lab-methods-1} --- Métodos de campo y de laboratorio I}

\begin{solution}{exo:b1:field-lab-methods-1:1}
A 1:10\,000, \qty{4.9}{m} se convierten en $4.9/10\,000 = \qty{0.49}{mm}$; a 1:25\,000,
$4.9/25\,000 = \qty{0.20}{mm}$, aproximadamente el grosor de un trazo de lápiz.
\end{solution}

\begin{solution}{exo:b1:field-lab-methods-1:2}
En Reikiavik $D = \ang{-11.28}$: $52 - 11.28 = \ang{040.7}$. En Wellington
$D = \ang{+23.23}$: $52 + 23.23 = \ang{075.2}$. La misma lectura de brújula significa
direcciones separadas \ang{34.5}.
\end{solution}

\begin{solution}{exo:b1:field-lab-methods-1:3}
$20\to015$. Con $t = \ang{15}$ y $p = \ang{20}$, el vector unitario hacia abajo es
\[
  (\cos\ang{20}\sin\ang{15},\ \cos\ang{20}\cos\ang{15},\ -\sin\ang{20}) =
  (0.2432,\ 0.9077,\ -0.3420).
\]
\end{solution}

\begin{solution}{exo:b1:field-lab-methods-1:4}
La media aritmética es $(350 + 358 + 6)/3 = \ang{238}$, casi opuesta a las tres lecturas. Los vectores
horizontales unitarios $(\sin a, \cos a)$ suman $(-0.1736 - 0.0349 + 0.1045,\ 0.9848 + 0.9994 + 0.9945) = (-0.1040, 2.9787)$, cuyo \omterm{def:b1:spherical-astronomy:horizontal}{acimut} es $\operatorname{atan2}(-0.1040, 2.9787) = \ang{-2.0}$, es decir,
\ang{358.0}: la media correcta, en el centro de las lecturas.
\end{solution}

\begin{solution}{exo:b1:field-lab-methods-1:5}
Media $175/5 = \ang{35.0}$. Desviaciones $-1, 2, 0, -2, 1$; suma de cuadrados 10;
$s = \sqrt{10/4} = \ang{1.58}$. Error típico $1.58/\sqrt5 = \ang{0.71}$.
\end{solution}

\begin{solution}{exo:b1:field-lab-methods-1:6}
Las normales son $(\sin\ang{30}, 0, \cos\ang{30}) = (0.5, 0, 0.866)$ y
$(-0.5, 0, 0.866)$; $\vec S = (0, 0, 1.732)$ es vertical: la media es un plano horizontal, y
$\bar R = 1.732/2 = 0.866$, cada lectura a \ang{30} de la media. Las dos lecturas no muestrean una
misma superficie con errores: pertenecen a dos flancos de un \omterm{def:b1:geological-maps:fold}{pliegue}, dos poblaciones distintas, y su
promedio no describe ninguna de ellas. Cada flanco se promedia por separado; juntos definen el eje del
\omterm{def:b1:geological-maps:fold}{pliegue}, la recta común a los dos planos, aquí horizontal y norte--sur.
\end{solution}

\begin{solution}{exo:b1:field-lab-methods-1:7}
$\tan p = \tan\ang{12}\cos(\ang{80} - \ang{120}) = 0.2126\times0.7660 =
0.1628$: $p = \ang{9.2}$, con las estrías inmersas hacia \ang{120}. Para la \omterm{def:b1:field-lab-methods-1:trend-plunge}{dirección} \ang{170},
$\cos(\ang{-90}) = 0$: la recta es horizontal, a lo largo del \omterm{def:b1:geological-maps:attitude}{rumbo}.
\end{solution}

\begin{solution}{exo:b1:field-lab-methods-1:8}
$0.47/0.20 = 2.35$ años después del 1 de enero de 2026, es decir, en la primavera de 2028. El 1 de enero
de 2030: $-0.47 + 4\times0.20 = \ang{+0.33}$, al este del norte geográfico.
\end{solution}

\begin{solution}{exo:b1:field-lab-methods-1:9}
Las tres normales están en el plano vertical de \omterm{def:b1:spherical-astronomy:horizontal}{acimut} \ang{080}, a \ang{10}, \ang{12} y \ang{14} de la
vertical; las dos exteriores son simétricas respecto de la central, así que su suma queda sobre ella: la
media es $080/12$, con $\theta = \ang{2}, 0, \ang{2}$. Entonces $\bar R = (1 +
2\cos\ang{2})/3 = (1 + 1.998782)/3 = 0.999594$, y $\sqrt{2(1 - \bar R)} =
\sqrt{0.000812} = \qty{0.02850}{rad} = \ang{1.633}$; exactamente,
$\theta_{\text{rms}} = \sqrt{(4 + 0 + 4)/3} = \ang{1.633}$.
\end{solution}

\begin{solution}{exo:b1:field-lab-methods-1:10}
La sección es un círculo de radio $\sqrt{r^2 - h^2}$. Su media es
$\frac1r\int_0^r\sqrt{r^2 - h^2}\,\mathrm dh = \frac1r\cdot\frac{\pi r^2}{4}
= \frac{\pi r}{4}$, ya que la integral es el área de un cuarto de disco. Un grano de cuarzo de
\qty{0.50}{mm} de diámetro aparece en promedio con $0.50\times\pi/4 = \qty{0.39}{mm}$ de diámetro: una
\omterm{def:b1:field-lab-methods-1:thin-section}{lámina delgada} subestima el tamaño de los granos.
\end{solution}

\begin{solution}{exo:b1:field-lab-methods-1:11}
La rotación de $D$ alrededor de la vertical,
$(x, y, z) \mapsto (x\cos D + y\sin D,\ -x\sin D + y\cos D,\ z)$, envía
$(\sin\delta\sin\alpha, \sin\delta\cos\alpha, \cos\delta)$ a
$(\sin\delta\sin(\alpha + D), \sin\delta\cos(\alpha + D), \cos\delta)$: suma $D$ a cada \omterm{def:b1:geological-maps:attitude}{dirección de
buzamiento}. Al ser lineal, envía la resultante $\vec S$ a la resultante de las normales giradas; conserva
las longitudes y la componente vertical, así que el \omterm{def:b1:geological-maps:attitude}{buzamiento} medio $\arccos(S_U/\|\vec S\|)$ no cambia,
y hace girar la parte horizontal de $\vec S$ en $D$, de modo que la \omterm{def:b1:geological-maps:attitude}{dirección de buzamiento} media aumenta en
$D$. Corregir cada lectura por la \omterm{def:b1:spherical-astronomy:equatorial}{declinación}, o solo su media, da el mismo resultado.
\end{solution}

\begin{solution}{exo:b1:field-lab-methods-1:12}
La dispersión a lo largo de la línea de máximo \omterm{def:b1:geological-maps:attitude}{buzamiento} es una dispersión del \omterm{def:b1:geological-maps:attitude}{buzamiento}: el error típico
es $2/\sqrt N \le 0.5$ para $N \ge 16$. A través de ella, un giro de la normal de \ang{2} cambia la \omterm{def:b1:geological-maps:attitude}{dirección de
buzamiento} en $2/\sin\ang{12} = \ang{9.62}$; el error típico de la \omterm{def:b1:geological-maps:attitude}{dirección de buzamiento} media es
$9.62/\sqrt N \le 1$ para $N \ge 92.5$, es decir, 93 lecturas. Un \omterm{def:b1:geological-maps:attitude}{buzamiento} suave se conoce pronto; su
dirección, no.
\end{solution}

\begin{solution}{pb:b1:field-lab-methods-1:1}
\textbf{1.} Direcciones de \omterm{def:b1:geological-maps:attitude}{buzamiento}: $967/12 = \ang{80.6}$; \omterm{def:b1:geological-maps:attitude}{buzamientos}: $145/12 =
\ang{12.1}$.

\textbf{2.} \ang{355} y \ang{005} promedian \ang{180}, la \omterm{def:b1:field-lab-methods-1:trend-plunge}{dirección} opuesta. Aquí todas las direcciones de
\omterm{def:b1:geological-maps:attitude}{buzamiento} están entre \ang{076} y \ang{085}, lejos del salto en \ang{000}, así que la media aritmética
funciona.

\textbf{3.} $\|\vec n\|^2 = \sin^2\delta(\sin^2\alpha + \cos^2\alpha) +
\cos^2\delta = 1$. La línea de máximo \omterm{def:b1:geological-maps:attitude}{buzamiento} hacia abajo tiene la \omterm{def:b1:field-lab-methods-1:trend-plunge}{dirección}
$(\cos\delta\sin\alpha, \cos\delta\cos\alpha, -\sin\delta)$, y su producto escalar con $\vec n$ es
$\sin\delta\cos\delta - \cos\delta\sin\delta = 0$; la línea de \omterm{def:b1:geological-maps:attitude}{rumbo} es horizontal y perpendicular a la
\omterm{def:b1:geological-maps:attitude}{dirección de buzamiento}, $(\cos\alpha, -\sin\alpha, 0)$, y su producto escalar con $\vec n$ es
$\sin\delta\sin\alpha\cos\alpha - \sin\delta\cos\alpha\sin\alpha = 0$.

\textbf{4.} La lectura 11 es la misma que la 2: $(0.2222, 0.0352, 0.9744)$. Lectura 12:
\[
  (\sin\ang{11}\sin\ang{82},\ \sin\ang{11}\cos\ang{82},\ \cos\ang{11}) =
  (0.1890,\ 0.0266,\ 0.9816).
\]

\textbf{5.} $\vec S = (2.4747,\ 0.4082,\ 11.7294)$.

\textbf{6.} $\|\vec S\| = \sqrt{6.1241 + 0.1666 + 137.5788} = 11.99457$, y
$\bar R = 11.99457/12 = 0.999547$.

\textbf{7.} $\bar\alpha = \operatorname{atan2}(2.4747, 0.4082) = \ang{80.6}$ y
$\bar\delta = \arccos(11.7294/11.99457) = \arccos 0.977893 = \ang{12.1}$: la media vectorial es
$080.6/12.1$, la misma que las medias aritméticas aquí.

\textbf{8.} Direcciones de \omterm{def:b1:geological-maps:attitude}{buzamiento}: $s = \ang{2.94}$; \omterm{def:b1:geological-maps:attitude}{buzamientos}: $s = \ang{1.68}$.

\textbf{9.} $\theta_{\text{rms}} \approx \sqrt{2\times0.000453} =
\qty{0.0301}{rad} = \ang{1.72}$.

\textbf{10.} $\sqrt{1.68^2 + (\sin\ang{12.1}\times2.94)^2} = \sqrt{2.82 +
0.38} = \ang{1.79}$, cerca de \ang{1.72}: las lecturas se comportan como dice el modelo de error, dentro de
los efectos del redondeo al grado y de una muestra pequeña.

\textbf{11.} El \omterm{def:b1:geological-maps:attitude}{buzamiento}: aporta 2.82 de los 3.20 (en grados cuadrados). La \omterm{def:b1:geological-maps:attitude}{dirección de buzamiento}
cuenta solo a través de $\sin\delta = 0.21$, porque en un estrato suave un cambio de \omterm{def:b1:geological-maps:attitude}{dirección de
buzamiento} gira muy poco el plano.

\textbf{12.} \omterm{def:b1:geological-maps:attitude}{Buzamiento} medio: $1.68/\sqrt{12} = \ang{0.48}$; \omterm{def:b1:geological-maps:attitude}{dirección de buzamiento} media:
$2.94/\sqrt{12} = \ang{0.85}$.

\textbf{13.} No: \ang{2.9} de \omterm{def:b1:geological-maps:attitude}{dirección de buzamiento} en un estrato que buza \ang{12} es solo
$2.9\sin\ang{12.1} = \ang{0.6}$ de giro del plano, menos que los \ang{1.7} de los \omterm{def:b1:geological-maps:attitude}{buzamientos}. Cualquier
pequeña ondulación de una superficie de estratificación suave dispersa mucho su \omterm{def:b1:geological-maps:attitude}{dirección de buzamiento}.

\textbf{14.} $D = \ang{-2.5}$: $83.1 - 2.5 = \ang{080.6}$; su media es $080.6/12.1$, la misma que la del
primer geólogo.

\textbf{15.} El \omterm{def:b1:geological-maps:attitude}{buzamiento} se lee con el clinómetro, frente a la vertical que da la gravedad; la aguja no
interviene.

\textbf{16.} Sí: una rotación alrededor de la vertical mueve cada lectura y la media en el mismo ángulo
(\cref{exo:b1:field-lab-methods-1:11}); la respuesta es la misma.

\textbf{17.} $10\times0.2 = \ang{2}$ de error en cada \omterm{def:b1:geological-maps:attitude}{dirección de buzamiento}, más del doble del error típico de
\ang{0.85}: un error sistemático que el promedio no puede detectar.

\textbf{18.} $80.6 - 90 = \ang{350.6}$; la \omterm{def:b1:geological-maps:attitude}{actitud} es \omterm{def:b1:geological-maps:attitude}{rumbo} \ang{351}, \omterm{def:b1:geological-maps:attitude}{buzamiento} \ang{12} hacia el este.

\textbf{19.} Direcciones de \omterm{def:b1:geological-maps:attitude}{buzamiento} $080.6$ y $080.6$: ninguna diferencia a esta precisión; \omterm{def:b1:geological-maps:attitude}{buzamientos}
$12.1$ y $12.2$: \ang{0.1}.

\textbf{20.} Las normales $(\sin\ang{12.2}\sin\ang{80.6},
\sin\ang{12.2}\cos\ang{80.6}, \cos\ang{12.2})$ y $\vec S/\|\vec S\|$ forman un ángulo de \ang{0.13}.

\textbf{21.} $1.72/\sqrt{12} = \ang{0.50}$; la diferencia, \ang{0.13}, es una cuarta parte de ese valor: no es
significativa. La brújula y los sondeos concuerdan.

\textbf{22.} Cuatro veces más: 48 lecturas.

\textbf{23.} Por ejemplo: un ajuste de declinación erróneo o desfasado; una brújula desviada por hierro
cercano (un vehículo, herramientas, rocas magnéticas); un clinómetro desajustado; lecturas tomadas solo en
las superficies mejor expuestas, que pueden no representar el flanco. El promedio reduce los errores
aleatorios, no los sistemáticos.

\textbf{24.} La \omterm{def:b1:geological-maps:attitude}{actitud} media del flanco oriental es \textbf{081/12}: \omterm{def:b1:geological-maps:attitude}{dirección de buzamiento} \ang{081},
\omterm{def:b1:geological-maps:attitude}{buzamiento} \ang{12}.
\end{solution}
